A validated, GPU-resident implementation of screened REBO2 was used to run an autonomous, hypothesis-driven study of how atomic-scale architecture controls the fracture of graphene-derived sheets, with the reference images serving as a source of design principles rather than templates. The findings, all statements about the \rebo{} model under athermal quasi-static uniaxial tension, are:
\begin{enumerate}
\item Strength is governed by the minimum load-bearing solid fraction across the load multiplied by a net-section strength that depends on ligament width and alignment (30--32\,N/m for straight, load-parallel ligaments wider than $\sim$20\,\AA; 17--19\,N/m for 10--13\,\AA{} round-pore ligaments; 12--20\,N/m for random networks), not on pore shape, porosity per se, hierarchy depth or scale ratio.
\item Load-path orientation is the largest lever: identical porosity yields 25 vs 4\,N/m for slits parallel vs perpendicular to the load, and 16 vs 12\,N/m for coarse veins along vs across the load.
\item Hierarchy adds no strength or toughness at matched mass unless its levels are aligned with the load; the only nested designs that beat the coarse-only limit combine load-elongated fine pores with veins (17.5\,N/m with square veins, progressive; 19.4\,N/m with load-parallel veins, stepwise; both 2.0\,J/m$^2$).
\item The fracture mode is a load-transfer (fibre-bundle) property: many parallel ligaments of moderate dispersion fail progressively and retain load, few equivalent wide strips, clustered defects and single flaws fail in one avalanche; strength dispersion alone does not create progressivity.
\item Cracks and holes larger than $\approx$20\,\AA{} sit on a lattice-trapping plateau (18--19.5\,N/m for 20--50\,\AA{} flaws), far above Griffith scaling.
\item Disorder, gradients and flaw halos are costs, not toughening mechanisms, in this model; apparent optima disappeared under seed replication, and the atomic registry of pore edges alone shifts fine-mesh strength by $\approx$12\%.
\item Twelve holdout designs predicted before simulation gave a median strength error of 14\% and the correct single-avalanche/multi-event classification in 10 of 12 cases; the misses identified a chirality-dependent ligament strength, the additivity of alignment effects across scales (best nested design 19.4\,N/m), the weakest-section rule for halos and gradients, and a new mechanism, en-echelon slit coalescence, at which load-path alignment is lost (8.9 instead of the predicted 23.9\,N/m).
\end{enumerate}
The platform, database (\nRuns{} records), structures, trajectories, figures, movies and this report are delivered as one reproducible archive; every simulation can be regenerated from its record.
